%%% SECTION 10: PATIENT INSTRUCTIONS FOR PHYSICAL AI ONCOLOGY TRIALS %%%

Patient Instructions: Physical AI Trials \cite{patient-instructions-paper}
provides clear, accessible, patient-facing documentation for participants in
Physical AI oncology trials. This document demonstrates that control has shifted
towards the patient's side, with comprehensive information about every robot
type participants may encounter. The patient instructions complement the
informed consent requirements of the adapted 21 CFR Part 50 \cite{cfr50-adapt}
and the patient rights provisions of AB 2847 from the Clinical Trial Site
Complete Documentation \cite{site-docs}.

\subsection{Purpose and Design Philosophy}
\label{subsec:instructions-purpose}

The patient instructions were designed following three principles: plain
language accessibility ensuring patients of all education levels can understand
the information, consistent structure for each robot type enabling easy
comparison and reference, and emphasis on patient rights and comfort
establishing that robots serve patients rather than the reverse.

Each robot type entry follows a standardized format: what the robot does and
why it is used in the trial, how the robot interacts with the patient (what
the patient will see, hear, and feel), safety features protecting the patient
during interaction, the patient's rights specific to that robot type (including
the right to refuse and request human alternatives), and what to do if the
patient has concerns or experiences discomfort.

This design philosophy reflects the National Platform's commitment to
patient-centered trial design. NCI information on clinical trial safety
\cite{nci-trials-safety} documents common patient concerns about trial
participation; Physical AI introduces additional concerns about robotic
interaction that these instructions specifically address. By providing
comprehensive, understandable information about each robot type, patients can
make informed decisions about their participation and exercise their rights
effectively.

\subsection{Surgical Robots}
\label{subsec:instructions-surgical}

Surgical robots (Robot Type 1) include the da Vinci system, Hugo RAS, and
Versius. Patient instructions explain that surgical robots assist the surgeon
in performing precise operations through small incisions. Patients learn that
the robot does not operate independently but follows the surgeon's commands
in real time. Safety features described include redundant control systems,
force limitation mechanisms, and the ability of the surgical team to stop the
robot instantly. Patient rights include the right to ask questions about the
robot before surgery, the right to request a non-robotic surgical approach,
and the right to have the surgical team explain each step during the procedure.
Cross-referencing with USL surgical robot scores \cite{usl-standard} confirms
that the da Vinci dVRK achieves the highest clinical trial readiness score
(Dimension D: 7.0) of any robot evaluated.

\subsection{Collaborative Robots (Cobots)}
\label{subsec:instructions-cobots}

Cobots (Robot Type 2) include the Franka Panda, Kinova Gen3, and xArm 7.
Patient instructions explain that cobots assist clinical staff with tasks
such as preparing medications, handling laboratory samples, and supporting
patient positioning. Patients learn that cobots are designed to work safely
alongside humans, with built-in sensors that detect unexpected contact and
automatically stop movement. Cobot interactions are typically indirect, with
the robot handling materials rather than touching the patient directly, though
some tasks (such as handing items to patients) involve proximity interaction.

\subsection{Radiotherapy Positioning Robots}
\label{subsec:instructions-rt-positioning}

Radiotherapy positioning robots (Robot Type 3) assist with precise patient
positioning for radiation therapy treatments. Patient instructions explain that
accurate positioning is critical for targeting tumors while protecting healthy
tissue, and that robotic positioning achieves sub-millimeter accuracy exceeding
manual positioning capability. Patients learn that the positioning process
involves the robot gently adjusting the treatment table or patient support
system, with the patient able to communicate with the therapy team throughout
the process. Safety features include continuous position monitoring, force
limits that prevent excessive pressure, and immediate stop capability.

\subsection{Needle-Placement Robots}
\label{subsec:instructions-needle}

Needle-placement robots (Robot Type 4) guide needles for biopsies and targeted
drug delivery with precision that exceeds manual capability. Patient
instructions explain that robotic guidance reduces the number of needle
attempts needed, shortens procedure time, and improves targeting accuracy.
Patients learn that the robot positions and guides the needle while a qualified
physician oversees the entire procedure. Comfort measures described include
local anesthesia, continuous communication with the clinical team, and the
ability to pause the procedure at any time.

\subsection{Social Companion Robots}
\label{subsec:instructions-companion}

Social companion robots (Robot Type 5) represent a unique aspect of Physical AI
oncology trials with no analog in traditional trials. Patient instructions
explain that companion robots provide emotional support, information delivery,
and well-being monitoring during trial participation. These robots can answer
questions about the trial process, provide relaxation exercises, offer
entertainment during treatment sessions, and alert clinical staff if the
patient appears distressed.

Patient instructions emphasize that companion robot interactions are entirely
optional and that patients may choose to interact with human staff exclusively.
The companion robot's ability to monitor patient mood and well-being is
explained transparently, with patients informed about what data is collected
and how it is used. This robot type addresses the emotional and psychological
needs of cancer patients, demonstrating that Physical AI is about patient
well-being and convenience as well as clinical efficiency.

\subsection{Humanoid Robots}
\label{subsec:instructions-humanoid}

Humanoid robots (Robot Type 6) including Tesla Optimus, Digit, and Atlas serve
multi-functional roles in trial sites including patient guidance through
facility areas, logistics support such as supply delivery and specimen
transport, and general assistance. Patient instructions address the unique
consideration of human-shaped robots: patients may find humanoid robots more
approachable due to their familiar form, but some patients may find them
unsettling. Instructions normalize both reactions and emphasize that patients
can request that humanoid robots not approach them.

\subsection{Radiotherapy Motion-Tracking Robots}
\label{subsec:instructions-rt-tracking}

Radiotherapy motion-tracking robots (Robot Type 7) continuously monitor patient
position during radiation therapy to ensure treatment accuracy. Patient
instructions explain that breathing and involuntary movement can shift the
treatment target, and that motion-tracking robots detect and compensate for
this movement in real time. Patients learn that the system may pause treatment
automatically if movement exceeds safe thresholds, and that this is a safety
feature rather than a malfunction. Safety interlocks ensure that radiation is
delivered only when positioning is confirmed accurate.

\subsection{Imaging Robots}
\label{subsec:instructions-imaging}

Imaging robots (Robot Type 8) perform diagnostic and monitoring imaging
including robotic CT positioning, MRI-compatible robotic systems, and
ultrasound guidance. Patient instructions explain that robotic imaging systems
can achieve more consistent positioning, faster scan times, and in some cases
reduced radiation exposure compared to manual imaging. DICOM \cite{dicom}
standards ensure that images are stored and shared in standardized formats
accessible to all members of the clinical team. Patients learn that imaging
data is protected under HIPAA \cite{hipaa} and the data protection provisions
of SB 892.

\subsection{Steerable Needle Robots}
\label{subsec:instructions-steerable}

Steerable needle robots (Robot Type 9) represent advanced guidance systems for
accessing complex anatomical targets. Unlike traditional straight-needle
insertion, steerable needle robots can navigate curved paths through tissue,
reaching targets that would be inaccessible or require multiple attempts with
manual techniques. Patient instructions explain the precision advantages and
note that the technology is particularly valuable for tumors in difficult
locations, reducing procedure time and improving targeting accuracy.

\subsection{Rehabilitation Exoskeletons}
\label{subsec:instructions-rehab}

Rehabilitation exoskeletons (Robot Type 10) support post-treatment recovery for
oncology patients. Patient instructions explain that exoskeletons provide
powered assistance for movement, helping patients rebuild strength and mobility
after surgery or prolonged treatment. The exoskeleton adapts to the patient's
capabilities, providing more support when needed and gradually reducing
assistance as the patient strengthens. Patients learn that the exoskeleton
monitors their movement patterns and reports progress to the clinical team.

\subsection{Patient Benefits Summary}
\label{subsec:instructions-benefits}

Physical AI oncology trials offer patients multiple documented benefits across
the ten robot types. Convenience benefits include reduced wait times through
automated scheduling, streamlined processes through robotic assistance, and
24-hour availability of robotic monitoring and support. Cost benefits include
lower treatment costs through reduced staffing overhead, fewer unnecessary
clinic visits through AI-guided monitoring, and shorter hospital stays through
precision surgical robotics. Quality benefits include sub-millimeter surgical
precision, continuous rather than periodic monitoring, AI-optimized treatment
plans, and comprehensive data-driven care. Emotional benefits include companion
robot support, reduced treatment anxiety through transparent information, and
empowerment through comprehensive patient rights.

NCI information on trial costs \cite{nci-trials-paying} provides context for
the financial benefits, while NCI safety information \cite{nci-trials-safety}
addresses patient concerns about trial participation. The patient instructions
implement the patient rights provisions of the adapted 21 CFR Part 50
\cite{cfr50-adapt}, the consent requirements of the adapted ICH E6(R3)
\cite{ich-adapt}, and the legally enforceable rights from AB 2847 in the
Clinical Trial Site Complete Documentation \cite{site-docs}, creating a
comprehensive patient communication framework that supports informed,
empowered participation in Physical AI oncology trials.

\subsection{Robot Type Summary Table}
\label{subsec:instructions-summary-table}

Table~\ref{tab:robot-summary} provides a comprehensive reference summarizing all
ten robot types with their primary clinical functions, representative
manufacturers, USL-relevant readiness information, and patient interaction
characteristics.

\begin{longtable}{c p{2.2cm} p{2.5cm} p{2.5cm} p{3cm} l}
\caption{Summary of All Ten Robot Types for Physical AI Oncology Trials}
\label{tab:robot-summary} \\
\toprule
\textbf{No.} & \textbf{Robot Type} & \textbf{Representative Models} & \textbf{Primary Function} & \textbf{Patient Interaction} & \textbf{Tier} \\
\midrule
\endfirsthead
\toprule
\textbf{No.} & \textbf{Robot Type} & \textbf{Representative Models} & \textbf{Primary Function} & \textbf{Patient Interaction} & \textbf{Tier} \\
\midrule
\endhead
\bottomrule
\endfoot
1 & Surgical & da Vinci, Hugo, Versius & Tumor resection, biopsy guidance & Direct invasive contact & 3 \\
2 & Cobots & Franka, Kinova, xArm & Lab automation, med prep & Indirect, material handling & 1-2 \\
3 & RT Positioning & Specialty systems & Radiation therapy alignment & Direct positioning contact & 2 \\
4 & Needle-Placement & Image-guided systems & Biopsy, drug delivery & Direct invasive contact & 3 \\
5 & Companion & Social robots & Emotional support, info & Non-contact interaction & 1 \\
6 & Humanoid & Optimus, Digit, Atlas & Logistics, guidance & Proximity, non-contact & 1 \\
7 & RT Tracking & Motion systems & Real-time motion tracking & Non-contact monitoring & 2 \\
8 & Imaging & Robotic scan systems & Diagnostic, monitoring & Positioning contact & 2 \\
9 & Steerable Needle & Advanced guidance & Complex target access & Direct invasive contact & 3 \\
10 & Rehab Exoskeleton & Wearable robots & Post-treatment recovery & Full-body wearable contact & 2 \\
\end{longtable}

The tier classification in the rightmost column directly determines the
regulatory requirements, consent elements, monitoring intensity, and operator
qualifications for each robot type as detailed in the adapted 21 CFR Part 50
\cite{cfr50-adapt} three-tier framework. Tier 3 robots (surgical, needle-
placement, steerable needle) require the most comprehensive informed consent
with procedure-specific elements, real-time multi-channel monitoring, and
qualified physician presence. Tier 2 robots (RT positioning, RT tracking,
imaging, rehabilitation) require enhanced consent and continuous telemetry.
Tier 1 robots (cobots in non-contact roles, companion, humanoid logistics)
require standard Physical AI disclosure with periodic monitoring.

\subsection{Benefits by Robot Category}
\label{subsec:instructions-category-benefits}

Each robot category offers distinct patient benefits that collectively create
a comprehensive care experience unavailable in traditional oncology trials.

Surgical robots provide precision benefits: sub-millimeter accuracy in tumor
targeting, reduced blood loss through minimally invasive approaches, smaller
incisions resulting in less post-operative pain, and shorter hospital stays.
The da Vinci platform's track record of 14 million procedures across 69
countries \cite{usl-standard} provides extensive safety data supporting
patient confidence.

Cobots provide efficiency benefits: faster medication preparation reducing
wait times, error-free sample labeling and tracking, consistent specimen
handling that preserves sample integrity, and around-the-clock availability
that eliminates delays from staff shift changes. The Franka Panda's highest
overall USL score (7.4) reflects the operational readiness that translates to
reliable patient service.

Companion robots provide psychological benefits that are unique to Physical AI
trials. Cancer patients experience anxiety, isolation, and information
overload. Companion robots provide on-demand emotional support, consistent and
patient responses to repeated questions, relaxation exercises tailored to
individual preferences, and a non-judgmental presence during difficult
treatment sessions. These benefits complement but do not replace human
emotional support from family, friends, and counselors.

Humanoid robots provide logistics benefits that improve the patient
experience: reliable supply delivery ensuring that treatment materials are
available when needed, specimen transport that maintains chain-of-custody
without human error, patient guidance through facility areas reducing
confusion, and environmental monitoring that maintains comfortable conditions.

Rehabilitation exoskeletons provide recovery benefits: personalized mobility
assistance that adapts to the patient's current capabilities, progressive
strength training with quantified progress tracking, motivation through
visible improvement metrics, and reduced dependency on physical therapy
staff for routine exercises.

Radiotherapy positioning robots provide accuracy benefits: sub-millimeter
patient alignment that exceeds manual positioning capability, consistent
reproducibility across multiple treatment sessions reducing inter-fraction
variation, faster setup times that reduce the duration patients must remain
immobile on the treatment table, and real-time position verification that
ensures radiation is delivered to the intended target. For patients undergoing
multi-week courses of radiation therapy, the cumulative benefit of faster and
more accurate positioning translates to reduced treatment fatigue, fewer
position-related treatment interruptions, and improved confidence in treatment
accuracy.

Needle-placement robots provide precision benefits: accurate first-attempt
needle positioning that reduces the number of insertion attempts, shorter
procedure duration resulting in reduced patient discomfort, improved targeting
accuracy for biopsies that increases diagnostic yield, and access to tumor
locations that would be difficult or impossible to reach with manual needle
insertion. For patients requiring repeated biopsies during long-term
monitoring, robotic needle placement reduces the cumulative discomfort and
anxiety associated with each procedure.

Radiotherapy motion-tracking robots provide safety benefits: continuous
position monitoring that ensures radiation is delivered only when the target
is accurately aligned, real-time compensation for breathing and involuntary
movement that allows treatment to proceed without requiring patients to hold
uncomfortable breath-hold positions, automatic treatment pauses when motion
exceeds safe thresholds that prevent radiation delivery to non-target tissue,
and reduced treatment margins enabled by improved motion management that
decreases radiation exposure to surrounding healthy tissue.

Imaging robots provide diagnostic benefits: more consistent positioning
across serial imaging studies that improves the reliability of treatment
response assessment, faster scan times that reduce the duration of
uncomfortable imaging positions, automated image quality optimization that
reduces the need for repeat scans, and AI-assisted image analysis that can
detect subtle changes in tumor size and characteristics earlier than manual
review. For patients in long-term monitoring (Stage 9 of the patient journey),
these benefits accumulate across multiple imaging sessions.

Steerable needle robots provide access benefits: the ability to navigate
curved paths through tissue reaching tumors in anatomically challenging
locations, reduced tissue damage compared to straight-needle approaches that
may require multiple insertion trajectories, shorter procedure times for
complex targeting scenarios, and the potential to reach targets that would
otherwise require open surgical approaches. These benefits are particularly
valuable for patients with tumors near critical structures where precision is
paramount for both treatment efficacy and avoidance of complications.

\subsection{Patient Rights Across Robot Types}
\label{subsec:instructions-rights-summary}

The patient rights established by AB 2847 \cite{site-docs} and implemented
through the adapted 21 CFR Part 50 \cite{cfr50-adapt} apply uniformly across
all ten robot types, ensuring that patients maintain consistent protections
regardless of which Physical AI systems are involved in their care. However,
the practical exercise of these rights varies by robot type and tier
classification.

The right to refuse robotic interaction has different implications depending on
the robot type. Refusing a Tier 1 robot (companion, humanoid logistics) has
minimal impact on clinical care because these systems support rather than
deliver treatment. Refusing a Tier 2 robot (positioning, imaging, tracking,
rehabilitation) may require alternative manual procedures that take longer but
achieve comparable outcomes. Refusing a Tier 3 robot (surgical, needle-
placement, steerable needle) requires the clinical team to provide a
conventional alternative (open surgery instead of robotic surgery, manual
needle placement instead of robotic guidance), which the trial protocol must
pre-specify as part of the Physical AI safety plan.

The right to real-time information during robotic procedures is implemented
differently across tiers. During Tier 1 interactions, real-time information is
available through companion robot interfaces or staff communication. During
Tier 2 interactions, the Physical AI operator provides continuous verbal
updates to the patient about the robot's actions and status. During Tier 3
interactions, the surgical team provides procedure updates through established
communication protocols, with the Physical AI console displaying relevant
status information visible to the operating team.

The right to a human advocate is exercised most frequently during Tier 3
procedures where patients may feel most vulnerable. The human advocate serves
as the patient's representative in the operating environment, monitoring the
procedure from the patient's perspective and authorized to request
clarification or raise concerns with the surgical team on the patient's behalf.
For Tier 1 and Tier 2 interactions, the patient may designate a family member,
friend, or patient navigator as their advocate, with access to the interaction
area subject to safety clearance requirements.

\subsection{Communication and Accessibility Standards}
\label{subsec:instructions-accessibility}

The patient instructions are designed to meet accessibility standards that
ensure all patients can understand and use the information effectively. The
document follows plain language guidelines targeting a sixth-grade reading
level for core content, with technical details available in supplementary
sections for patients who desire more information. Medical terminology is
defined at first use and collected in a patient glossary at the end of the
document.

Visual design standards include consistent formatting across all ten robot type
sections, enabling patients to quickly locate information about specific robots.
Each section uses the same sequence of information (what the robot does, how it
interacts with you, safety features, your rights, what to do if concerned),
creating a predictable structure that patients can navigate independently.

Language accessibility requires that the patient instructions be available in
all languages spoken by more than five percent of the patient population at
each trial site. Translation must be performed by qualified medical translators
with review by native-speaking clinical staff to ensure accuracy of both
medical content and cultural appropriateness. Cultural adaptation extends
beyond translation to address culturally specific concerns about robotic
interaction, privacy expectations, and healthcare decision-making practices.

Accessibility for patients with disabilities follows ADA and Section 508
standards. Large-print versions are available for patients with visual
impairments. Audio versions are available for patients who cannot read printed
materials. Sign language interpretation is available for consent discussions
with deaf patients. The electronic version of the patient instructions is
compatible with screen readers and other assistive technologies.

The patient instructions are updated whenever Physical AI systems at a trial
site change, new robot types are added, or consent requirements are modified.
Updated versions are distributed to all active participants along with a
summary of changes, supporting the ongoing consent process required by the
adapted Part 50. Version tracking ensures that every participant's record
documents which version of the patient instructions they received and when
they received it, creating a complete information delivery audit trail.
